% Created 2026-05-25 Mon 15:17
% Intended LaTeX compiler: pdflatex
\documentclass[11pt]{article}
\usepackage[utf8]{inputenc}
\usepackage[T1]{fontenc}
\usepackage{graphicx}
\usepackage{longtable}
\usepackage{wrapfig}
\usepackage{rotating}
\usepackage[normalem]{ulem}
\usepackage{amsmath}
\usepackage{amssymb}
\usepackage{capt-of}
\usepackage{hyperref}
\author{MeCantSpell}
\date{\today}
\title{Responsive Web Design Certificate Notes}
\hypersetup{
 pdfauthor={MeCantSpell},
 pdftitle={Responsive Web Design Certificate Notes},
 pdfkeywords={},
 pdfsubject={},
 pdfcreator={Emacs 30.2 (Org mode 9.7.11)}, 
 pdflang={English}}
\begin{document}

\maketitle
\tableofcontents

This file contains the notes for my Responsive Web Design Certificate. I am following a Bad Website
Club bootcamp!
\section{Computers Basics}
\label{sec:org3253279}
Computers are machines that allow people to accomplish certain tasks including playing games,
browsing the internet and even writing this note.
As a developer it's important to take a look at the parts of a computer as it's the machine for
which we will be developing software for.
\subsection{Parts of a Computer}
\label{sec:orgd06402a}
\textbf{Motherboard} : The motherboard is the main part of a computer as it's the circuit board that houses
all of the other important components. It can be thought of as the skull of the computer.
\textbf{CPU} : CPU stands for Central Processing Unit, it's needed for
the computer to function. It's a small square chip tht goes into a CPU Socket on the
motherboard. It's speed is calculated in MHz and GHz (Million and Billion Cycles per Second
respectively). It's responsible for calculations, providing CPU time for tasks and other functions
as well. It can be thought of as the brain of the computer.
\textbf{RAM} : RAM stands for Random Access Memory and and it's needed to store important details required
to perform functions by the CPU. The CPU takes data and information stored in the RAM when it needs
to complete its tasks. Memory stored in the RAM is not persistent among
restarts of the computer. Any important data must therefore be saved to the permanent memory. It can
be thought of as the short term memory of the computer.
\textbf{HDD/SSD} : HDD stands for Hard Disk Drive and is one of the most important storage devices in the
history of computers. Utilizing platters and heads it stores data across the platters for
retrieval. While in today's modern computers they are overshadowed by their newer counterpart of
SSD, standing for Solid State Drives, they are still cheaper and can be found in often larger
sizes. These storage devices are persistent and can be considered as the long term memory of the
computer.
\textbf{PSU} : PSU stands for Power Supply Unit and is the device by which a computer recieves electricity
to function. It converts the power from out powerlines into power suitable to be consumed by the
computer(AC --> DC). Without it the computer will not function. The PSU can be thought of as the
digestive tract of the computer.
\textbf{Expansion Cards} : In order to give the computer additional functionality, the motherboard will
house certain expansion slots in which expansion cards can be slot in. Expansion cards that are
commonly used include:
\begin{enumerate}
\item Video Cards - Also known as GPUs which help with rendering video graphics.
\item Sound Cards - Responsible for the audio output of a computer
\item Network Cards - Responsible for connecting the computer to the internet.
\end{enumerate}
\subsection{Working Effectively With Input Devices}
\label{sec:org956f5c7}
There are certain ways in which we can improve the efficiency of our work on the computer with
computer peripherals. We'll take a look at some of the tips we can incorporate in our daily tech
lives.
\begin{enumerate}
\item Gripping the mouse gently in a relaxed manner
\item Using ergonomic peripherals
\item Keeping the mouse at the same height as the keyboard
\item Maintaining good posture \textbf{(No slouching!)}
\item Using keyboard shortcuts
\item \textbf{MOST IMPORTANTLY} Taking regular breaks
\end{enumerate}
\subsection{Internet Service Providers}
\label{sec:org3bf179b}
Internet Service Providers or ISPs are companies that sell people the access to the global
internet. There are 3 types of ISPs.
\begin{enumerate}
\item \textbf{Tier 1 ISPs} : These are ISPs that can handle most of their network traffic independently. They
can also access most of the internet without any restrictions.
E.g.: Lumen Technology (Rank 1)
\item \textbf{Tier 2 ISPs} : These are the ISPs that handle the network of a country mostly. They can access
some parts of the internet without restrictions but have to pay to access some parts of the
internet.
E.g.: Vodafone (Purchases access from Arelion)
\item \textbf{Tier 3 ISPs} : These are the ISPs that handle regional networks and solely rely on larger ISPs
to access parts of the internet.
E.g.: Hutch (4 million users though)
\end{enumerate}

These ISPs will offer different types of connections to access the internet. We'll take a look at
some connections here.
\begin{enumerate}
\item \textbf{Fibre Connections} : Transmits data via glass or plastic fibre cables, allows for very high speed
connections and data exchanges. (We have fibre UwU)
\item \textbf{Cable Connections} : Transmits data via television cables making them readily available and
accessible in many regions.
\item \textbf{DSL Connections} : Transmits data via landline cables making them much more accessible than
cable connections.
\item \textbf{Dial-UP Connections} : They also transmit data via landline cables, however they are only able
to do one thing at a time (either calls or internet). This is an archaic form of internet
connection.
\item \textbf{Satellite Connections} : They transmit data between an array of satellites.
\item \textbf{Cellular Network Connections (4G/5G)} : They transmit data between cellular
\end{enumerate}
\subsection{Sign In Security}
\label{sec:org8d3c5b8}
When it comes to your computer it's important to make sure it's password protected. When setting up
a password use a password that is long and complex, and not an easy password. Passwords should also
not use personal information such as names, dates and ages. On your computer, if your device allows
it, setting up biometric methods (fingerprints, face recognition) to sign in is considered a safer
approach when compared to traditional passwords.
When setting up your passwords for services, using 2FA (two factor authorization) and a password
manager (E.g.:\href{https://proton.me/pass}{Proton Pass}) are recommended.
\subsection{Developer Tools}
\label{sec:orge899b0f}
In order to create software, developers rely on a bunch of tools to make their work more productive
and efficient.
\begin{enumerate}
\item \textbf{Computers} : Developers have to use a computer in order to write their code. However compared to
the average user, a developer would require a computer with higher processing power and RAM so
they can carryout resource intensive tasks.
\item \textbf{IDEs/Code Editors} : The software that a developer chooses to write their code in is referred to
as a code editor. The software in which this note is being written on (Emacs) is also a code
editor. IDEs on the other hand stand for Integrated Development Environments and provide powerful
features including code completion, debugging and integrated terminal support (All of this can be
kind of added to Emacs, so I dunno\ldots{} it's also an IDE or something).
\item \textbf{VC Systems} : When writing code, code is likely to break. Also, when a team is involved in
development, different developers will code in different styles and ways. A VC (Version Control)
system allows developers to retrieve previous versions of a codebase in case something goes
wrong. Think of it as an elaborate undo button for code. The most popular version control system
is GitHub (we'll get magit later, don't worry)
\item \textbf{Package Managers} : When writing complex code, developers will make use of packages and
libraries. Package managers simplify the process of adding, updating and removing these packages.
E.g.: PIP for Python
\item \textbf{Testing Frameworks} : In order to test whether the code that developers have written is
functioning as expected. While you can manually test your code by running it through test cases
and documenting the pass/fail of your code, a testing framework allows you to automatically do
that.
E.g.: pytest for Python
\item \textbf{Web Browsers} : Developers also need to test what the code looks like visually beyond just the
functionality. This can be done using browsers like Mozilla Firefox, Vivaldi, Opera GX,
Thorium. Testing on different browsers is important to make sure the product is accessible by as
many users as possible.
\end{enumerate}
\subsection{File Management Applications}
\label{sec:orgb1e179f}
A file management application makes it easy to store, organize and view the files on your
computer. On Windows the default file management application is File Explorer whereas on Mac it's
finder.
File Explorer can be accessed via the \uline{Start Menu by pressing the Windows Key on your keyboard or by}
\uline{clicking the Windows Icon on your taskbar.} Using File Explorer you can pin, copy and move files
around.
These file management applications provide features like tags,search and Smart Folders.
\subsection{Naming Files for Web Applications}
\label{sec:org0053a3c}
HTML files used for web applications have the extension \textbf{.html} after the name of the file. For web
applications we will also use files like CSS files and JavaScript files which have the extensions
\textbf{.css} and \textbf{.js} respectively.
While there are no set conventions for naming your web application files, it is recommended that you
name them using sensible names and avoid using special characters that can make the name look
confusing. Importantly, when working with a team it is vital to follow your team's style-guide when
naming you files.
The exception to this is \textbf{index.html} which is a special file that represents the main page of the
website. It's the first file that is loaded when you visit a website.
\subsection{File/Folder Organization for Web Applications}
\label{sec:org26758c8}
When building out web applications (or any applications in general), it makes sense to have all the
files and folders you need arranged in a way that makes your workflow productive.
In a file system a root directory is referred to as the top-level directory in a file system, all
other files and sub-directories located within this root directory are referenced via the root
directory.
Think of the root directory as the District you live in, and the sub directories point to the street
you live in and the files within them finally point to your house (I am in your walls).
A root directory can be represented by a single dot \textbf{(.)}
You can have sub directories within your directory for the assets used in your website, CSS
stylesheets, JS scripts and perhaps one for your webpages.
A best practice directory will look like this:
\begin{verbatim}
.
|___/assets
|___/css
|___/pages
|___/js
|___/admin
|___/blogs
|___index.html
|___README.org
\end{verbatim}
Personally I like to add a little TODO.org so that I have an idea what I have to do next and a
little note of what I have already done, or failed doing.
\subsection{Creating, Moving and Deleting Files}
\label{sec:orgf3a55d8}
In windows while in File Explorer, once you have entered your desire directory you can create a new
file \uline{by right-clicking and going to the "New" menu and selecting the type of file you would like to}
\uline{create.} You can also rename files \uline{by either selecting your file and clicking the rename button on}
\uline{the top, right-clicking and choosing the rename option, selecting your file and hitting F2 or by}
\uline{double-clicking on the name underneath the file.} By default you don't have to enter the extension as
files will be saved in their appropriate file extension, in order to change the file extension you
can either \uline{enable file extensions on your systema and manually edit the extensions of your desired}
\uline{files or by editing the extension name in your "Save As" prompt when you are initially saving.}
Deleting files in Windows is really easy, you just have to \uline{select the files you want to delete and}
\uline{press the del key, use the delete option on the top or right-click and select the delete option.}
In order to move files you will either have to copy/cut and paste them elsewhere. You can copy files
\uline{by selecting the desired files and right-clicking to copy or using the Ctrl-C shortcut}. You can cut
files \uline{by selecting the desired files and right-clicking to cut orusing the Ctrl-X shortcut.} Finally,
to paste these files \uline{simply right-click and select paste or use the Ctrl-V shortcut.}
In order to select multiple files on Windows you can either \uline{click and drag your mouse or use}
\uline{Ctrl+click over the relevant files (Ctrl+click really helps when the files aren't adjacent).}
\subsection{Searching for Files}
\label{sec:orgb3dc6d7}
In Windows you used to be able to search for the files you needed using the taskbar, this no longer
works properly. You could go to the directory where you know the file is and search from there. But honestly
if you have internalized all the details from the previous chapter you won't need to search for
anything as you will be well organized to always know where everything is.
\subsection{File Types used in Web Applications}
\label{sec:orgac59670}
Earlier in \ref{sec:org0053a3c} we looked at files like .html, .css and .js that we
will use when we develop web applications. Apart from these we will also use some other files types
such as:
\begin{enumerate}
\item \textbf{JPEG} : Joint Photographic Experts Group
\item \textbf{PNG} : Portable Network Graphics
\item \textbf{GIF} : Graphics Interchange Format (pronounced JIF)
\item \textbf{SVG} : Scalable Vector Graphics
\end{enumerate}
Which are some image formats used in web applications.
\begin{enumerate}
\setcounter{enumi}{4}
\item \textbf{MP3} : A lossy audio format
\item \textbf{WAV} : A lossless audio format
\item \textbf{MP4} : The most common video format
\item \textbf{WebM} : A high-quality lossless open-source video format
\item \textbf{TFF} : A font format
\item \textbf{WOFF} : A modern font format
\item \textbf{WOFF2} : WOFF but 2
\item \textbf{ZIP} : Compressed folders
\item \textbf{Markdown} : Files that can be used to create structured documents (use org though, org files can
be exported as markdown, html and even pdf)
\end{enumerate}
\subsection{Web Browsers}
\label{sec:org67d543b}
Web browsers are appilcations that allow you to access the internet. The current most popular
browsers are Microsoft Edge, Mozilla Firefox, Google Chrome and Safari. Your system should come with
a default browser and if you feel like you want to change your browser it's just a matter of looking
up your preferred browser on your defualt browser and downloading the one you want.
On some systems you can install a browser using your operating system's package manager like on
MacOS with homebrew or Arch Linux.
\subsection{Web Terms}
\label{sec:org188d990}
We've already looked at what a web browser is, now let's take a look at some terms that are often
confused for each other, like search engine and website.
\textbf{Web Browsers} : Web browsers are the applications used to navigate the world wide web, they are
commonly found with features such as an address bar, rendering engine, bookmarks, extensions and
various security pages.
\textbf{Websites} : Websites on the other hand are the content that is made available on the World Wide Web
and can be accessed through your web browser. They are a collection of web pages and other files,
and consist of a domain name and the web pages that make them. Websites often include multimedia and
can also be static (content rarely changes) or dynamic (content frequently changes) depending on the
needs of the publisher.
\textbf{Search Engines} : Search engines are web-based tools designed to help users find the content that
they are looking for on the internet. Key aspects of search-engines include web crawlers also known
as spiders or bots, which are prorgams that systematically browse the web to discover and index new
content. They also have complex algorithms that rank the content that they have indexed and display
them to the user based on factors such as relevance, authority, popularity and user experience.
Popular search engines include Google, Bing, Yahoo, and DuckDuckGo.
Broswers, search engines and websites are all interconnected. It can be simply thought of this,
websites are the content of the internet, they can be accessed via search engines or browsers. If
you know the website you want to go to you don't need a search engine, however if you are not sure
of where to go you can use a search engine to help you get there.
\subsection{Advanced Search Engine Capabilities}
\label{sec:orgc1f7c73}
When using your search engine there are certain ways that you can search more efficiently. For
example if you want to search for a specific phrase you can put that phrase in between \textbf{quotation
marks ("")}. If you want to search for specific words but they do not have to necessarily be together
you can use the \textbf{plus symbol (+)} before the words. If you want your search engine to not return
search results with certain words you can instead use the \textbf{minus symbol (-)}. To search within a
specific website or domain you can use the following syntax \textbf{site:}
A few example searches may look like this:
\begin{verbatim}
"freecodecamp curriculum"
\end{verbatim}
This one will return only results that have the exact phrase of "freecodecamp curriculum".
\begin{verbatim}
site:freecodecamp.org/news +python -curriculum
\end{verbatim}
This one looks at any webpages with the domain freecodecamp.org and within their "news" path.
Note that search engines are not case-sensitive and essentially convert all searches to lower
case. The above capabilities will be available on Google Search Engine, and also work on
Startpage. However, it makes sense to read the documentation on the search engine you are using to
see the exact syntax that is used for these advanced functionalities.
For more tips on using the Google Search Engine, and by extension Startpage, visit this \href{https://woodward.libguides.com/google/basicsearch}{site}.
\section{HTML}
\label{sec:org24399aa}
HTML stands for Hypertext Markup Language and is a language used for building webpages. When you
visit a website, its content including headings, paragraphs, sentences, images, videos and links are
all made of HTML.
While HTML is enough to create a functional website, and it is considered a best-practice to have
website functionality in HTML. A modern day website also needs CSS for styling and JavaScript for
interactivity.
\subsection{Elements}
\label{sec:orgf0b2d26}
HTML is made up of elements. This elements usually have an opening tag and an ending tag. Elements
that do not have ending tags are referred to as void elements. These elements can also be classified
as semantic and non-semantic elements. Simply put, semantic elements are elements that have meaning
whereas non-semantic elements are elements that do not have meaning.
Using semantic elements over non-semantic elements has certain benefits including increased
accessibility for screen readers, enhanced SEO for search engine indexing and finally to increase
code readability.
After the opening tag of an element we will enter in the content. If the element also has a closing
tag, once we are done entering our content we will use the closing tag to close it, otherwise if it
is a void element there is no need.
HTML tags are case-insensitive, however the widely adopted best-practice is to use simple-case
letters when typing them in. Most tags will start with \textbf{<} and end with \textbf{>}. For example:
\begin{verbatim}
<p>
\end{verbatim}
A closing tag however will also include a forward slash \textbf{/} as well. Taking the above example:
\begin{verbatim}
</p>
\end{verbatim}
We have explained that void tags do not have a closing tag earlier, and therefore they will follow
the convention of the opening tags. However, one can write a void tag as such and it will be valid:
\begin{verbatim}
<img/>
\end{verbatim}
Many code formatters indeed choose to format void tags like this. However the HTML spec states that
the presence of \textbf{/} "does not mark the start tag as self-closing but instead is unnecessary and has no
effect of any kind".
Some elements will take in values known as attributes that define the way the content will behave in
said tag is displayed. The basic syntax for an attribute is the attribute name followed by an \textbf{=} sign
For example the image element will take in attributes such as \textbf{src} which
indicates the source/location of the image and \textbf{alt} which gives alternative text for the user in case
the image isn't displayed correctly. For example:
\begin{verbatim}
<img src="https://cdn.freecodecamp.org/curriculum/cat-photo-app/catsjpg" alt="Two tabby kittens sleeping together on a couch." >
\end{verbatim}
Another example is the attributes used for hyperlinks. Hyperlinks are created using the anchor
element \textbf{<a>}, and attributes such as \textbf{href} and \textbf{target} can be used. "href" refers to href stands for
hypertext reference and is the reference link of the hyperlink and the target attribute indicates
how the hyperlink will open up. For example:
\begin{verbatim}
<a href="https://www.freecodecamp.org" target="_blank"> Visit freeCodeCamp </a>
\end{verbatim}
The \textbf{\_blank} value of the target attribute indicates that this hyperlink must be opened in a new tab.
Just like how some elements are unique in their syntax some attributes are also unique in their
syntax like the \textbf{checked} attribute which indicates whether or not an input element of the type
checkbox is checked by default.
\begin{verbatim}
<input type="checkbox" checked/>
\end{verbatim}
Such attributes are referred to as boolean attributes because just like boolean they can either be
present or absent. Other boolean attributes include disabled, readonly and required which describe
the state of an element.
\subsection{Comments}
\label{sec:org4ba9f6c}
In order to type in comments in HTML we start them with the characters \textbf{<!--} and end them with the
characters \textbf{-->}. A full comment will look something like this:
\begin{verbatim}
<!--This is an HTML comment.-->
\end{verbatim}
\subsection{Boilerplate}
\label{sec:org6dfd1ec}
A boilerplate is a readymade template for making websites using HTML. The boilerplate can be
considered as the skeleton of the site, and you can add the bones and meat to it using the specific
HTML that is required for the instance, as well as using CSS and JS.
A boilerplate is important because it ensures that your pages are structured correctly and work well
across multiple web browsers. Using a boilerplate helps you avoid common mistakes and follow
best-practices.
As you personalize your own boilerplate, you will realize that it saves time when creating new
projects and it will become your default starting point for any web project.
\subsection{UTF-8}
\label{sec:org21a225e}
UTF-8 stands for UCS Transformation Format-8, UCS in turn stands for Universal Character Set. It is
the standardized character encoding widely used on the web (99\% of the web uses UTF-8). The other
unicode standard that exists (existed) is UTF-16 which is now nearly obsolete with only 0.004\% of
cwebpages using it.
Using UTF-8 and by that extent, setting it as the characterset in a meta element in your webpages,
is crucial to maintaining the standards of the internet as it supports every character in the
Unicode set from all writing systems, languages and techincal symbols. It's why you can type in
Tamil and Sinhala and Emojis.
The syntax that is therefore required to enter into your webpage withing the <head></head> element
is:
\begin{verbatim}
<meta charset="UTF-8" />
\end{verbatim}
Having this code in your boilerplate will also prove efficient.
\subsection{HTML Entities}
\label{sec:orgd3ca5d3}
As you may have seen when creating your web document you will have used a varying amount of syntax
including symbols like "<",">" and "=". However sometimes, there may be a need to display those
symbols on your webpage. An HTML entity is therefore referred to as a set of characters used to
represent a reserved charcter in HTML. They are also known as character references.
There are three different types of character references but it is important to note that all of them
start with an \textbf{ampersand(\&)} and end with a \textbf{semi-colon(;)} in common.
\begin{enumerate}
\item \textbf{Named character references}: Named character references start with an \textbf{ampersand(\&)} and end with a
\textbf{semi-colon(;)}. These are references that are used for actual HTML elements and prevents the
parser from confusing them for the actual elements.
\item \textbf{Decimal numeric reference}: Here the reference starts with the usual \textbf{ampersand(\&)} but is followed
immediately by a \textbf{hash(\#)} symbol. Here specific decimal numerics that correlate to the relevant
symbol are used instead. Then it is finally ended with a \textbf{semi-colon(;)}
\item \textbf{Hexadecimal numeric reference}: Just like the decimal numeric reference, it starts with an
\textbf{ampersand(\&)}, \textbf{hash(\#)} but is the followed by the simple letter \textbf{x} then it is followed by the ASCII
hexadecimal numerics that correlate to the relevant symbol. Finally it is also closed with a
\textbf{semi-colon(;)}.
\end{enumerate}
\subsection{Media Assets}
\label{sec:org729e860}
In order to enhance the visual aspect of a webpage or to pass out visual inforamtion we will use
media assets such as images, audios and videos. However it is important to note that using such
assets will take a heavier toll on the data consumed to retrieve them, loading times and
storage. Therefore using these media assets in an optimal way is encouraged and a best
practice. There are three important things to consider when using an image, they are the display
size, file format and compression.
When it comes to the display size, it is important to remember
that our media assets should match the size we are displaying on our webpage. If our webpage is
displaying an image that is 640x480 but the image has not been scaled down and is instead a
1920x1080 image, we are making the users download additional data and storing unnecessary data as
well.
When it comes to format it is important to use the most ideal format, while PNGs and JPEGs
are the most common formats for images, they are not the most ideal formats and instead you should
use WEBP or AVIF formats. Below we will give a list of optimal file formats for various media
assets:
\begin{enumerate}
\item \textbf{Images}: \textbf{WEBP}, AVIF, \textbf{SVG}
\item \textbf{Audios}: \textbf{OGG}, \textbf{WEBM}, WAV, M4A, MP3
\item \textbf{Videos}: \textbf{WEBM}, MP4
Note however that MP4 with codecs like H.264 and HVEC may open you up to legal issues
\end{enumerate}
Finally it must be noted that in order to make our files smaller and more efficient, we can make use
of compression algorithms. Note that compression takes two forms: lossy and lossless. Their names
are self-explanatory. One of them preserves the quality of the original, whereas the other loses
data to reduce its size.
Another important point to consider when using media assets in your webpages, is the license they
hold. By default all creative works released are released with all rights reserved, meaning the
creator or publisher holds all the rights for the media asset. Such assets cannot be used unless one
of three things has occurred:
\begin{enumerate}
\item You have expressed written permission by the author.
\item You have purchased a licensing fee.
\item You incorporate the media under fair use.
Fair use requires that your use of the original media is both limited and transformative.
\end{enumerate}
\subsection{Paths}
\label{sec:org922c1d8}
The goal of creating websites is to share content. Now, all of this content may not be on the same
webpage or it may not be feasible to have it all on the same webpage and therefore it is vital that
we find a way to like our content to wherever the user is currently. We have already seen the anchor
element that is used to link content and allow users to move to another site. We have also looked at
other elements such as images, videos, audios and iframes that allow us to link content of those
specific nature from a different source.
Let's take a look at the specificities of those sources. There are two types of paths that exist:
\begin{enumerate}
\item \textbf{Relative Paths}
These are paths that specify the location of a file relevant to the current file
\item \textbf{Absolute Paths}
These are paths that specify the location of a file directly.
\end{enumerate}
Think of it like this, let's say you are friends with 2 guys: Sam and Jack. Let's also assume that
your house is next to Jack's. If Sam is at your house and asks where Jack's house is, you are not
going to describe every street, corner and road you have to take to get there, you'd simply say
"Oh. He's house is on the right.", however if Sam asked you where Jack's house is from school, you
might have to explain how to get all the way to Jack's house from Sam's house, or the school or the
most common \textbf{root} you both share.
That's how these paths work. It's also important to understand that absolute paths and absolute URLs
are different as well. The former gives the absolute direction to the file within a file system
whereas the latter gives the absolute directions to the file including the protocols and for web
resources the domain name as well.
There are different instances of when we should use relative and absolute paths. For example when
testing on our own PC we might know the exact paths that the files need to access in order to
complete their tasks, but when we ship our program to another PC they might not have installed the
files in the same way and therefore using an absolute path there will cause the program to not
function correctly. Likewise, when we need to get an exact resource for say, an iframe, from another
webpage we cannot use a relative path and we have to give the absolute URL.
Knowing when to use and not to use absolute and relative paths is very important in web development.
When using relative paths there are some syntax that we should remember. Let us create a new example
directory and take a look through it:
\begin{verbatim}
Project/
|___Resources/
|___|___Images/
|___|___Videos/
|___Site/
|___|___index.html
|___|___styles.css
|___|___scripts.js
\end{verbatim}
Let's say our index.html needed to access our styles.css file. In order to for it to access it, we
will use \textbf{a single dot(.)} in order to indicate that we want to access a file that is in the
subdirectory that we are sharing. So in order to link the relevant file using relative path we will
use the following link:
\begin{verbatim}
./styles.css
\end{verbatim}
Let's say that there exists a file called favicon.svg in our images folder and we need to access it
via our index.html. In this case we will use \textbf{two double dots(..)} in order to access the parent
directory (i.e. Project/)
\begin{verbatim}
../Resources/Images/favicon.svg
\end{verbatim}
\subsection{Accessibility}
\label{sec:orgdcdd4f7}

\subsection{HTML Syntax}
\label{sec:orge315cbc}
This list has been created using the syntax that I have personally encountered and syntax that I am
much likely to use. For a full reference on all the HTML syntax please refer to this \href{https://developer.mozilla.org/en-US/docs/Web/HTML/Reference/Global\_attributes}{link}.
\begin{enumerate}
\item Anchor/Hyperlink - <a>,</a>
\begin{enumerate}
\item Hypertext Reference - href
\item Browsing Context - target
("\textsubscript{self}" for the same tab, "\textsubscript{blank}" for a new tab or new window, there are other values as
well but since they are values that deal with the now obsolete frames we will ignore them)
\end{enumerate}
\item Blink - <blink>,</blink>
Blink is now a depecrated HTML syntax
\item Body - <body>,</body>
\item Bold - <b>,</b>
\item Content Division - <div>,</div>
\item Figure - <figure>,</figure>
The figure is a new element introduced in HTML5 with the goal of adding semantics to figures
including images, video clips, drawings and other visual content. Within the figure element you
can add <figcaption></figcaption> that will add a caption to the image/s as well as the images
themselves.
\item Head - <head>,</head>
\item Headings - <h1 - 6>,</h1 - 6>
\begin{enumerate}
\item Create a unique id for CSS styling - id
\item Create a class to share CSS styling among multiple elements - class
\end{enumerate}
\item HTML - <html>,</html>
Webpage language - lang
\item HTML initialization - <!DOCTYPE html>
\item Image - <img/>
\begin{enumerate}
\item Source - src
\item Alternative Text - alt
\end{enumerate}
\item Input - <input/>
\begin{enumerate}
\item Type - type
Sets the type of input obtained (checkbox, password, email, url)
\item Autocrrect - autocorrect
\item Defualt check for checkbox - checked
\item Disable element by default - disabled
\item Make element immutable - readonly
\item Makes the element mandatorily accept a value - required
\end{enumerate}
\item Italicize - <i>,</i>
\item Linking to external resources - <link/>
\begin{enumerate}
\item Defining the relationship - rel
Sets the relationship type (stylesheet, preconnect, icon)
\item Hyperlink reference - href
\end{enumerate}
\item Main - <main>,</main>
\item Metadata - <meta/>
Defining the character set - charset
Defining a metadata - name
Entering content for the metadata - content
\item Moving around - <marquee>,</marquee>
Marquee is now a depecrated HTML syntax
\item Paragraph - <p>,</p>
\begin{enumerate}
\item Create a unique id for CSS styling - id
\item Create a class to share CSS styling among multiple elements - class
\end{enumerate}
\item Semantic Sections - <section>,</section>
\item Title - <title>,</title>
Title is displayed on the browser tab
\begin{enumerate}
\item Create a unique id for CSS styling - id
\item Create a class to share CSS styling among multiple elements - class
\end{enumerate}
\item Emphasis (Semantic Italicization) - <em>,</em>
\item Strong (Semantic Bold) - <strong>,</strong>
\item Adding scripting languages - <script>,</script>
\item Content Division - <div>,</div>
Division elements are non-semantic meaning they don't carry meaning. They are useful for
grouping together elements that need to share the same CSS styling. However, more often than not
there is usually a more appropriate semantic element that can be used instead.
\begin{enumerate}
\item Create a unique id for CSS styling - id
\item Create a class to share CSS styling among multiple elements - class
\end{enumerate}
\item Section - <section>,</section>
\item Audio - <audio>,</audio>
\begin{enumerate}
\item Source - src
\item Audio Controls - controls
\item Autoplay Audio - autoplay
\item Mute audio by default - muted
\item Loop audio - loop
\item Video Format - type
\end{enumerate}
\item Source - <source/>
Used as a nested element to store multiple sources for its parent element.
\item Video - <video>,</video>       
\begin{enumerate}
\item Source - src
\item Video Controls - controls
\item Autoplay video - autoplay
\item Mute audio by default - muted
\item Loop video - loop
\item Adjust width - width
\item Put a video thumbnail while the video is downloading - poster
\item Video Format - type
\end{enumerate}
\item Audio and video text tracks - <track/>
\begin{enumerate}
\item Source - src
\item Kind/Type of track - kind
Values can be captions (includes descriptions of sounds), subtitles, chapters (for
navigation), descriptions (summary of the media) and metadata
\item Language - srclang
\item Language Label - label
\item Enable the track by default - default
\end{enumerate}
\item Inline Frame used to embed content such as videos, maps and documents - <iframe>,</iframe>
\begin{enumerate}
\item Source - src
\item HTML Source - srcdoc
\item Width - width
\item Height - height
\item Title which works like alt-text - title
\item Allowing what the iframe can do - allow
\end{enumerate}
\item SVG - <svg>,</svg>
\begin{enumerate}
\item Width - width
\item Height - height
\item How much the browser will display - viewbox
\end{enumerate}
\item SVG Circles - <circle/>
\begin{enumerate}
\item Border - stroke
\item Border width - stroke-width
\item Colour fill - fill
\end{enumerate}
\item SVG Paths - <path/>
\begin{enumerate}
\item Path colour - stroke
\item Path width - stroke-width
\item Area between path colour - fill
\end{enumerate}
\item Embedded content - <embed>,</embed>
General purpose embedding tool with little to no use
\item Embedded content - <object>,</object>    
General purpose embedding tool with little to no use
\end{enumerate}
\section{CSS}
\label{sec:orge23c915}

\subsection{CSS Syntax}
\label{sec:org5d14f81}
\begin{enumerate}
\item Anchor Elements
\begin{enumerate}
\item Default link state - :link
\item Visited links - :visited
\item Links being hovered upon - :hover
\item Links being focused upon - :focus
I personally cannot tell the difference between :focus and :active
\item Links that are being actively interacted with - :active
\end{enumerate}
When creating a CSS style for anchor elements it is important to style them in this specific
order:
\begin{verbatim}
:link, :visited, :hover, :focus, :active
\end{verbatim}
\end{enumerate}
\end{document}
